\section{valid/ready 握手与背压}

同步互连最小事务可以用 valid/ready 描述：发送者在 valid 为 1 时给出稳定负载，接收者在 ready 为 1 时表示本拍能够接收；只有二者在同一时钟边沿都为 1，事务才真正转移。若 ready 拉低，发送者必须保持 valid 和负载不变，不能悄悄换成下一笔。

\begin{longtable}{L{0.18\linewidth}L{0.28\linewidth}L{0.44\linewidth}}
\toprule
valid & ready & 边沿发生的事 \\
\midrule
0 & 0/1 & 没有有效事务 \\
1 & 0 & 背压；发送者保持全部负载 \\
0 & 1 & 接收者空闲等待 \\
1 & 1 & 当前事务完成一次转移 \\
\bottomrule
\end{longtable}

组合地把 ready 反接到 valid，容易在多个模块串联后形成长组合路径甚至组合环。高频设计会插入 skid buffer 或寄存器切片，在不丢事务的前提下切断这条路径。

\section{AXI/APB 的通道与排序边界}

APB 面向低带宽寄存器设备，一次传输用 setup 和 access 阶段完成；AXI 把读地址、读数据、写地址、写数据、写响应分成独立通道，允许多笔事务在途。独立通道意味着写地址与写数据可以不同拍到达，从设备必须分别接收并在配对后执行，不能假设两者同时出现。

事务 ID 用来区分在途请求。互连可以让不同 ID 的响应交错返回，但同一 ID 或带顺序要求的事务必须遵守协议定义。超时、decode error 和 slave error 也要沿响应通道返回；未映射地址不能让总线永远等待。

\begin{keyidea}
地址译码回答“交给谁”，握手回答“何时接收”，ID 与顺序规则回答“哪个响应属于哪笔请求”，错误响应回答“失败怎样闭合”。四者缺一都不是完整互连。
\end{keyidea}
